\begin{lemma}
\label{editorial-source-lemma-geometrically-irreducible-base-change-transcendental}
Let $K/k$ be a field extension. The following are equivalent
\begin{enumerate}
\item $K$ is geometrically irreducible over $k$, and
\item the induced extension $K(t)/k(t)$ of purely transcendental extensions
is geometrically irreducible.
\end{enumerate}
\end{lemma}

\begin{proof}
Assume (1). Let $\Omega$ be an
algebraic closure of the original
$k(t)$. The algebra $K\otimes_k\Omega$
has irreducible spectrum by (1).
Put $U=k[t]\setminus\{0\}$ and
$V=K[t]\setminus\{0\}$, keeping
the same indeterminate $t$.
There is an isomorphism
$$
K\otimes_k k(t)\longrightarrow U^{-1}K[t],
\qquad a\otimes(P(t)/Q(t))\longmapsto aP(t)/Q(t).
$$
The inverse sends
$(\sum_i a_it^i)/Q(t)$ to
$\sum_i a_i\otimes(t^i/Q(t))$.
The quotient relations for the
original denominators verify
independence of representatives,
and the formulas give both composites.
Localizing further at the images
of all of $V$ gives precisely the
original fraction field $K(t)$.

The two original scalar maps
$$
(K\otimes_k k(t))\otimes_{k(t)}\Omega
\longrightarrow K\otimes_k\Omega,
\qquad(a\otimes b)\otimes c\longmapsto a\otimes bc
$$
and $a\otimes c\mapsto(a\otimes1)\otimes c$
are inverse. Localization commutes
with these tensors: the fraction
$b/v\otimes c$ goes to
$(b\otimes c)/(v\otimes1)$, with
inverse determined by inverting
the same original $v$. Hence
$K(t)\otimes_{k(t)}\Omega$ is a
localization of $K\otimes_k\Omega$
at exactly those images of $V$.
It is nonzero because it is a
tensor product of two nonzero
fields over $k(t)$; extending
$1$ to vector-space bases proves
this directly. A nonzero
localization of a ring with
prime nilradical again has prime
nilradical, as proved in Lemma
\ref{editorial-source-lemma-subalgebra-geometrically-irreducible}.
Thus its spectrum is irreducible.
The algebraic-closure test of Lemma
\ref{editorial-source-lemma-geometrically-irreducible}
proves (2).

Conversely, assume (2) and retain
any original extension $k'/k$.
Put $A=K\otimes_k k'$ and let $W$
be the multiplicative subset of
$A[t]$ generated by the images of
all nonzero polynomials of $K[t]$
and of $k'[t]$. Both coefficient
maps into $A$ are injective:
extension of a vector-space basis
over $k$ proves this. More strongly,
the leading coefficient of each
such polynomial is a unit of $A$,
with inverse given by the inverse
of the original nonzero field
coefficient in the same tensor
factor. Multiplication by a
polynomial with unit leading
coefficient is injective on $A[t]$:
if $H$ is nonzero of degree $n$
and leading coefficient $h_n$,
the leading coefficient of its
product with a degree-$d$
polynomial of leading coefficient
$c$ is exactly $c h_n\neq0$.
Products of these injective
multiplication maps remain
injective. Thus every element
of $W$ is a nonzerodivisor and
$A\to A[t]\to W^{-1}A[t]$ is
injective, without assuming $A$
reduced or a domain.

The full receiving algebra is
exactly
$$
K(t)\otimes_{k(t)}k'(t)\simeq W^{-1}A[t].
$$
To construct the map from left
to right, use the original
coefficient embeddings and send
$$
\frac{F(t)}{f(t)}\otimes\frac{G(t)}{g(t)}
\longmapsto\frac{F(t)G(t)}{f(t)g(t)},
\qquad
f\in K[t]\setminus\{0\},\quad
g\in k'[t]\setminus\{0\}.
$$
The field-fraction relations hold
after precisely these denominators
are inverted. The two images of
every rational function in $k(t)$
agree, so the map is balanced over
$k(t)$. For the inverse, send
$a\otimes b\in A$ to $a\otimes b$
in the left-hand tensor and send
$t$ to $t\otimes1=1\otimes t$.
The image of each original
$f\in K[t]\setminus\{0\}$ is
the unit $f(t)\otimes1$, and
each original nonzero $g\in k'[t]$
gives the unit $1\otimes g(t)$.
The universal property therefore
extends this map to $W^{-1}A[t]$.
The composites fix $A$, $t$ and
each specified inverse denominator,
so they are identities on their
generating sets. In particular,
the injection just proved is
the original tensor-constant map
used in the source argument.

By (2), this receiving tensor
has nonempty irreducible spectrum,
so its nilradical is a proper
prime. Its inverse image under
the injection from $A$ is exactly
$\operatorname{Nil}(A)$, because
an injection reflects the vanishing
of every power. Hence $A$ is
nonzero with prime nilradical,
and its spectrum is irreducible.
This works for every original
$k'/k$, proving (1).

Iteration gives the same equivalence
after adjoining any specified finite
list of common independent
indeterminates. At each step the
proved one-variable equivalence
applies to the actual preceding
field extension and retains all
its polynomial denominators.
The empty list is the identity
extension, so it also satisfies
the assertion.
\end{proof}